% !PRINTED 182
\weberpage{182}
%<<CONT
mehr vorausgesetzt wird, dass die Determinante der Function $\varphi$
verschwindet. Diese Discriminante ist
\weq{1}{\psi(x) = [\Sigma x_i \varphi'(x'_i)]^2 - 4\varphi(x)\varphi(x').}
Setzen wir in ihr $x_i + \lambda x'_i$ an Stelle von $x_i$, so ergiebt sich,
wenn man f\"ur $\varphi(x + \lambda x')$ den Ausdruck \S.\,56, (11) substituirt
und $\Sigma x'_i \varphi'(x'_i) = 2\varphi(x')$ setzt,
\[
\psi(x + \lambda x') = \psi(x),
\]
woraus hervorgeht, dass $\psi(x)$ nur von $m - 1$ linearen Verbindungen
 der $m$ Ver\"anderlichen abh\"angt. Setzt man f\"ur $x'_1, x'_2 \ldots x'_m$
beliebige specielle Werthe, f\"ur die $\varphi(x')$ nicht verschwindet, und
nimmt $x'_1$ von Null verschieden an, so wird
% [sic] Im Druck steht x_n; die Variablen laufen sonst bis x_m.
\[
x_1'^2 \psi(x_1, x_2, \ldots x_n) = \psi(0, x_2 x'_1 - x_1 x'_2, \ldots x_m x'_1 - x_1 x'_m),
\]
also eine Function der $m - 1$ Ver\"anderlichen
\[
y_2 = x_2 x'_1 - x_1 x'_2, \ldots y_m = x_m x'_1 - x_1 x'_m.
\]
Man kann dies auch dadurch best\"atigen, dass
\[
\psi'(x_1), \ \psi'(x_2) \ldots \psi'(x_m)
\]
f\"ur $x_1, x_2 \ldots x_m = x'_1, x'_2 \ldots x'_m$ alle zugleich verschwinden,
dass also die Determinante der Function $\psi$ verschwinden muss.

Die Formel (1) enth\"alt nun ein sehr wichtiges Resultat, das
sich am besten \"ubersehen l\"asst, wenn wir dieser Formel folgende
Gestalt geben.

Wir setzen
\weq{2}{\begin{aligned}
\psi(x) &= -4\varphi(x')\chi(y_2, \ldots y_m) \\
\Sigma x_i \varphi'(x'_i) &= 2\sqrt{\varphi(x')}\, Y_1,
\end{aligned}}
und erhalten
\weq{3}{\varphi(x) = Y_1^2 + \chi(y_2, y_3 \ldots y_m),}
wodurch $\varphi(x)$ dargestellt ist als Summe aus einem Quadrat einer
% [sic] Im Druck steht hier das kleine y_1, während (2) und (3) Y_1 verwenden.
linearen Function $y_1$ und einer Function $\chi$ von $m - 1$ Variablen.

Daraus aber folgt der wichtige Satz

\gesperrt{Eine quadratische Form von $m$ Ver\"anderlichen l\"asst
sich immer darstellen als eine Summe von Quadraten
von $m$ oder weniger linearen Functionen.}

Dieser Satz ist offenbar richtig f\"ur eine Function von einer
Variablen, die ja selbst ein Quadrat ist. Aus (3) folgt aber die
Richtigkeit des Satzes f\"ur eine Function von $m$ Ver\"anderlichen,
wenn seine Richtigkeit f\"ur eine Function von $m - 1$ Ver\"anderlichen
 vorausgesetzt wird.
